\section{Conclusion}
Ce rapport a appliqué la démarche d'urbanisation au périmètre de la carte nationale d'identité électronique. L'étude de l'existant a montré un portail qui a réussi la généralisation du rendez-vous en ligne et une plateforme d'identité numérique souveraine déjà adoptée, mais un processus encore coupé en deux entre le numérique et le guichet, un contrôle des pièces manuel et tardif, et des canaux cloisonnés. La captation de la stratégie a traduit Digital Morocco 2030, la loi 55-19 et les priorités de la DGSN en cinq objectifs mesurables et cinq principes d'urbanisation.

Le système cible décline ces objectifs de manière cohérente sur quatre couches~: un processus P0 reconstruit autour de l'identité numérique, du contrôle documentaire automatique et de l'enrôlement allégé~; un plan d'occupation des sols de trente-deux blocs fonctionnels dont douze nouveaux~; quinze composants applicatifs orchestrés par un moteur de processus et reliés par une passerelle d'API~; une architecture technique en zones, hébergée souverainement, avec une plateforme d'inférence dédiée. Le prototype Bonita du sous-processus de pré-demande a confirmé la faisabilité de l'orchestration et des interfaces. La trajectoire en trois paliers permet de livrer de la valeur dès 2027 tout en maîtrisant les risques humains, techniques et réglementaires.

Deux perspectives prolongent ce travail. La première est l'extension de la démarche aux autres téléservices de la DGSN, pour lesquels le socle d'intégration, le moteur de processus et la plateforme d'IA proposés ici sont directement réutilisables. La seconde est la mesure réelle des gains~: les indicateurs définis dans la cartographie métier devront être instrumentés dès le premier palier afin d'ajuster les seuils de décision des modèles et de documenter, chiffres à l'appui, l'apport de l'intelligence artificielle au service public de l'identité.
